

\documentclass{article}
\usepackage[T1]{fontenc}
\usepackage[utf8]{inputenc}
\usepackage{lmodern}
\usepackage{tikz}
\usetikzlibrary{positioning,arrows.meta,shapes.geometric,fit}
\usepackage{amsmath}
\usepackage{amssymb}
\usepackage{adjustbox} % for rotation+scaling
\usepackage[a4paper, left=0.6in, right=0.6in, top=1in, bottom=1in]{geometry}
\pagestyle{empty}

\begin{document}

\textbf{(a) Between-class variance and failure condition.}

\medskip

\noindent\textbf{Setup.} Background: fraction $1-p$, $X_B\sim\mathrm{Unif}[-b,b]$, mean $\mu_B=0$.
Foreground: fraction $p$, $X_F\sim\mathrm{Unif}[1-a,1+a]$, mean $\mu_F=1$.
Assume no overlap: $b<1-a$.
For threshold $k$ let class~0 be $\{x\le k\}$ and class~1 be $\{x>k\}$.
Define
\[
\omega_0(k)=\Pr\{x\le k\},\quad \omega_1(k)=1-\omega_0(k),
\]
\[
\mu_0(k)=\mathbb{E}[x\mid x\le k],\quad \mu_1(k)=\mathbb{E}[x\mid x>k].
\]
Otsu maximizes
\[
\sigma_B^2(k)=\omega_0(k)\,\omega_1(k)\big(\mu_0(k)-\mu_1(k)\big)^2.
\]

\medskip

\noindent\textbf{Case 1: $k$ in the gap $(b,\,1-a)$.}
Then all background pixels are $\le k$ and all foreground pixels are $>k$, so
\[
\omega_0=1-p,\quad \omega_1=p,\quad \mu_0=0,\quad \mu_1=1,
\]
and
\[
\boxed{\sigma_B^2(k)=p(1-p)\qquad (k\in(b,\,1-a)).}
\]

\medskip

\noindent\textbf{Case 2: $k\in[-b,b]$ (inside background).}
Let
\[
\alpha(k)=\Pr_B\{X\le k\}=\frac{k+b}{2b}\in[0,1],
\]
then
\[
\omega_0(k)=(1-p)\alpha(k),\qquad \mu_0(k)=\frac{-b+k}{2}.
\]
With total mean $\mu_T=(1-p)\cdot0+p\cdot1=p$ and using
$\mu_1=(\mu_T-\omega_0\mu_0)/\omega_1$ we get the compact form
\[
\boxed{\sigma_B^2(k)=\frac{\omega_0(k)\big(\mu_0(k)-p\big)^2}{1-\omega_0(k)}
\qquad (k\in[-b,b]).}
\]

\medskip

\noindent\textbf{Failure condition.}
Compare $\max_{k\in[-b,b]}\sigma_B^2(k)$ with the gap value $p(1-p)$.
If the interior maximum exceeds $p(1-p)$, Otsu will choose a threshold inside
the background and thus fail to separate true classes.

Differentiate $\sigma_B^2$ (equivalently work in $x=\mu_0=(k-b)/2$). Evaluating the
derivative at the endpoint $k=b$ (i.e. $x=0$) gives
\[
\left.\frac{d\sigma_B^2}{dx}\right|_{x=0}=\frac{(2b-1)(p-1)}{b}.
\]
Since $p-1<0$, the derivative points toward larger $\sigma_B^2$ into $(-b,b)$
iff $2b-1>0$, i.e. iff $b>\tfrac12$. Therefore, under the no-overlap assumption
$b<1-a$, Otsu \emph{can} fail whenever
\[
\boxed{b>\tfrac12.}
\]
Intuition: a wide background (half-width $b>1/2$) can be split to produce a
larger between-class variance than the correct gap split $p(1-p)$.

\newpage

\textbf{(b) One iteration of $k$-means initialized at means $0$ and $1$.}

\medskip

\noindent With current means $\mu^{(0)}_1=0,\ \mu^{(0)}_2=1$ the assignment
boundary is the midpoint $x=0.5$.

\begin{itemize}
\item If $b\le\tfrac12$: the entire background $[-b,b]$ lies to the left of $0.5$,
  so all background points are assigned to cluster~1 and all foreground points
  to cluster~2. The recomputed means remain $(0,1)$ and $k$-means is stable
  after one iteration (correct separation).

\item If $b>\tfrac12$: a fraction of background, namely the interval $(0.5,b]$,
  lies to the right of $0.5$ and will be assigned to cluster~2 together with all
  foreground. The new cluster~2 mass is
  \[
  \omega_2' \;=\; p \;+\; (1-p)\frac{b-0.5}{2b},
  \]
  and its new mean is
  \[
  \mu_2'=\frac{p\cdot 1 \;+\; (1-p)\frac{b-0.5}{2b}\cdot\frac{0.5+b}{2}}
                {\,p + (1-p)\frac{b-0.5}{2b}\,}\;<1.
  \]
  Similarly $\mu_1'>0$. Hence both means move inward, the between-class distance
  shrinks, and the algorithm has mixed true classes (wrong clustering).
\end{itemize}

\medskip

\noindent\textbf{Agreement with (a).} This matches part (a): the critical value
is $b=\tfrac12$. If $b\le\tfrac12$ the midpoint $0.5$ cleanly separates the supports
and both Otsu and $k$-means succeed; if $b>\tfrac12$ the midpoint cuts into the
background, causing misassignment for $k$-means and producing a larger
between-class variance inside the background for Otsu. The root cause is the
large background spread relative to the inter-class distance $1$.



\end{document}


